\section{Lugar de ejecución de la propuesta}

El desarrollo del presente proyecto se llevará a cabo principalmente en la Universidad Nacional de Colombia, sede Palmira, en el marco del programa de Maestría en Ingeniería Agroindustrial, donde se realizarán las actividades relacionadas con el diseño, modelado, análisis y evaluación del sistema de secado propuesto.

De manera complementaria, la Universidad del Valle sede Melendez en Cali participará como institución aliada, aportando el desarrollo previo de la tecnología y brindando soporte técnico en el proceso de optimización y escalamiento del sistema.

Asimismo, se contempla la articulación con Cenicafé, ubicado en el municipio de Chinchiná, Caldas, como entidad de acompañamiento para la validación técnica y el análisis del desempeño del sistema de secado en el contexto del sector cafetero colombiano.

Adicionalmente, el proyecto incluye su proyección hacia un escenario de aplicación real en una planta de beneficio de café proyectada en el municipio de Sevilla, Valle del Cauca, donde se evaluará la viabilidad de implementación de la tecnología en condiciones operativas representativas.

Esta articulación entre el entorno académico, el acompañamiento investigativo y el contexto productivo permite garantizar la pertinencia, aplicabilidad y potencial de transferencia tecnológica de los resultados del proyecto.